% Part: computability
% Chapter: computability-theory
% Section: reducibility

\documentclass[../../../include/open-logic-section]{subfiles}

\begin{document}

\olfileid{cmp}{thy}{red}
\olsection{తగ్గించదగినత}

\begin{explain}
గణనీయంగా లెక్కించదగినదే కాని గణనీయం కాని కనీసం ఒక సమితి $K_0$
ఉందని ఇప్పుడు మనకు తెలుసు. అలాంటి ఇతర సమితులు కూడా ఉన్నాయని స్పష్టమే.
మళ్లీ మళ్లీ మూల సూత్రాల నుంచే మొదలుపెట్టకుండా, ఇతర సమితులకు ఈ ధర్మాలు
ఉన్నాయని చూపడానికి తగ్గించదగినత పద్ధతి శక్తివంతమైన సాధనాన్ని ఇస్తుంది.

సాధారణంగా, సమితి $A$ను సమితి~$B$కు ``తగ్గించడం'' అంటే,
\tetoken{మూలకాలు}{!!{element}s}~$B$లో ఉన్నాయా లేవా అనే సమాధానాలను,
అవే \tetoken{మూలకాలు}{!!{element}s}~$A$లో ఉన్నాయా లేవా అనే
సమాధానాలుగా మార్చే పద్ధతి. మనం ``అనేకం-ఒకటి తగ్గించదగినత'' అనే
భావనపై దృష్టి పెడతాం; అయితే భిన్న ధర్మాలు గల మరెన్నో తగ్గించదగినతా
భావనలు ఉన్నాయి. సామర్థ్య సమస్యలనూ పరిగణించాల్సిన గణనా సంక్లిష్టత
అధ్యయనంలో కూడా తగ్గించదగినతా భావనలు కేంద్రస్థానంలో ఉంటాయి. ఉదాహరణకు,
ఒక సమితి NPలో ఉండి, ప్రతి NP సమస్యనూ దానికి తగ్గించగలిగితే దాన్ని
``NP-సంపూర్ణం'' అంటారు. ఇక్కడ వాడే తగ్గింపు భావన కింద ఇచ్చేదానికి
పోలినదే; అదనంగా ఆ తగ్గింపును బహుపది కాలంలో గణించగలగాలి.

తగ్గింపు భావనను మనం ఇప్పటికే పరోక్షంగా వాడాం. సమితి~$K$ను ఇలా
నిర్వచించండి:
\[
K = \Setabs{x}{\cfind{x}(x) \fdefined},
\]
అంటే $K=\Setabs{x}{x\in W_x}$. ఆగే సమస్య అనిర్ణయనీయమని మనం ఇచ్చిన
నిరూపణ (\olref[hlt]{thm:halting-problem}) అత్యంత నేరుగా $K$ గణనీయం
కాదని చూపుతుంది.
\sourcecorrection{OLTECOMTHY-010}{ఆంగ్ల మూలం ఈ పేరా మొదటిలో ``notion of reduction notion'' అని ``notion''ను రెండుసార్లు రాసింది; పునరుక్తిని తొలగించి ఉద్దేశించిన వాక్యాన్ని అనువదించాం.}
\sourcecorrection{OLTECOMTHY-011}{ఆగే సమస్య గురించి ఆంగ్ల మూలం ``is unsolvable''కు బదులు ``in unsolvable'' అనే అక్షరదోషంతో రాసింది. ముందున్న సిద్ధాంత సూచనకు అనుగుణంగా సరైన అర్థాన్ని ఇచ్చాం.}
ఇప్పుడు $K_0$ అనేది
\[
K_0 = \Setabs{\tuple{e, x}}{\cfind{e}(x) \fdefined},
\]
అంటే $K_0=\Setabs{\tuple{e,x}}{x\in W_e}$. $K$ గణనీయం కాదనే ఏ
నిరూపణనైనా $K_0$ గణనీయం కాదనే నిరూపణగా విస్తరించడం సులభం:
$K_0$ గణనీయమైతే, $\tuple{x,x}$ అనేది $K_0$లో ఉందా అని అడగడం ద్వారా
\tetoken{ఒక మూలకం}{!!a{element}}~$x$ అనేది $K$లో ఉందా లేదా
నిర్ణయించవచ్చు. $x$ను $\tuple{x,x}$కు
పంపే ప్రమేయం~$f$, $K$ను $K_0$కు తగ్గించడానికి ఒక ఉదాహరణ.
\sourcecorrection{OLTECOMTHY-012}{ఆంగ్ల మూలం ముందుగా $K_0$ను $\tuple{e,x}$ క్రమంతో నిర్వచించినప్పటికీ, సమానమైన $W_e$ లక్షణీకరణలో జతను $\tuple{x,e}$గా తారుమారు చేసింది. సాధారణ జత-సంకేతీకరణ సమమితం కానందున, నిర్వచనానికి అనుగుణంగా $\tuple{e,x}$గా సరిచేశాం.}
\end{explain}

\begin{defn}
$A$, $B$లు సహజ సంఖ్యల సమితులు అనుకుందాం. ప్రతి సహజ సంఖ్య~$x$కు
\[
x \in A \quad \text{అయితేనే} \quad f(x) \in B.
\]
అయ్యే గణనీయ ప్రమేయం $f\colon\Nat\to\Nat$ను $A$ నుంచి~$B$కు
\emph{అనేకం-ఒకటి తగ్గింపు} అంటాం. అలాంటి తగ్గింపు~$f$ ఉంటే, $A$
అనేది~$B$కు \emph{అనేకం-ఒకటిగా తగ్గించదగినది} అని చెప్పి
$A\leq_m B$ అని రాస్తాం. $A$ అనేది $B$కు, అలాగే విలోమంగా కూడా
అనేకం-ఒకటిగా తగ్గించదగినవైతే, $A$, $B$లను \emph{అనేకం-ఒకటి
తుల్యాలు} అని చెప్పి $A\equiv_m B$ అని రాస్తాం.
\end{defn}

పై నిర్వచనంలోని ప్రమేయం~$f$ అంతఃక్షేపకమైతే, $A$ అనేది~$B$కు
\emph{ఒకటి-ఒకటిగా తగ్గించదగినది} అంటాం. కింద వర్ణించే తగ్గింపుల్లో
చాలావరకు ఈ బలమైన ఆవశ్యకత కూడా నెరవేరుతుంది; కానీ మనం ఆ వాస్తవాన్ని
వాడం.

\begin{digress}
ఒకటి-ఒకటి తగ్గించదగినత నిజంగానే అనేకం-ఒకటి తగ్గించదగినత కంటే బలమైన
ఆవశ్యకత; అయితే ఇది ఎంతమాత్రమూ స్పష్టమైన వాస్తవం కాదు. మరో మాటలో,
$A$ అనేది $B$కు అనేకం-ఒకటిగా తగ్గించదగినదే కాని $B$కు ఒకటి-ఒకటిగా
తగ్గించదగినది కాని అనంత సమితులు $A$,~$B$ ఉన్నాయి.
\end{digress}

\end{document}
